% ===========================================================================
\chapter{Invent\'ario de objetos e relat\'orio de malha}\label{ap:objetos}
% ===========================================================================

Esta se{\c{c}}\~ao reproduz os totais e o envelope do relat\'orio de malha emitido pelo build.
A tabela detalhada, com uma linha por objeto e as contagens de v\'ertices, faces, quads,
tri\^angulos, n-gons, arestas n\~ao variedade e arestas de borda, est\'a no arquivo
\texttt{docs/renders/mesh\_report.txt}, que \'e gerado automaticamente e versionado.

\section{Totais consolidados}

\begin{table}[htbp]
\centering
\caption{Totais de malha da cena (malha base, sem subdivis\~ao)}
\label{tab:totais-ap}
\footnotesize
\begin{tabular}{>{\raggedright\arraybackslash}p{7.4cm}r}
\toprule
\textbf{Grandeza} & \textbf{Valor} \\
\midrule
Objetos de malha & 130 \\
Faces & 32.632 \\
\quad quads & 29.272 \\
\quad tri\^angulos & 3.360 \\
\quad n-gons & 0 \\
Arestas n\~ao variedade & 0 \\
Arestas de borda & 1.776 \\
\bottomrule
\end{tabular}
\fonte{\texttt{docs/renders/mesh\_report.txt} (2026).}
\end{table}

\section{Envelope dimensional}

\begin{table}[htbp]
\centering
\caption{Envelope medido e alvo}
\label{tab:envelope-ap}
\footnotesize
\begin{tabular}{>{\raggedright\arraybackslash}p{4.6cm}rr>{\raggedright\arraybackslash}p{4.6cm}}
\toprule
\textbf{Grandeza} & \textbf{Medido (m)} & \textbf{Alvo (m)} & \textbf{Observa{\c{c}}\~ao} \\
\midrule
Comprimento (X) & 3,736 & 3,726 & $+1,0$\,cm \\
Largura (Y) & 1,718 & 1,718 & exato \\
Altura (Z) & 1,117 & 1,117 & exato \\
Solo (Z m\'inimo) & 0,0000 & 0,000 & exato \\
Centro em Y & 0,0000 & 0,000 & exato \\
\bottomrule
\end{tabular}
\fonte{\texttt{docs/renders/mesh\_report.txt} (2026).}
\end{table}

\section{Invent\'ario por grupo funcional}

\begin{table}[htbp]
\centering
\caption{Objetos de malha por grupo funcional}
\label{tab:grupos-ap}
\footnotesize
\begin{tabular}{>{\raggedright\arraybackslash}p{7.0cm}r}
\toprule
\textbf{Grupo} & \textbf{Objetos} \\
\midrule
Rodas (4 conjuntos de 13 pe{\c{c}}as) & 52 \\
Ventila{\c{c}}\~oes (\emph{louvers}, sa\'idas, duto) & 23 \\
Far\'ois e indicadores & 12 \\
Lanternas traseiras & 8 \\
Interior & 7 \\
Grelha frontal & 6 \\
Aerodin\^amica & 6 \\
Far\'ois auxiliares & 4 \\
Espelhos & 4 \\
Casco e cabine & 3 \\
Emblemas & 3 \\
Refletores do difusor & 2 \\
\midrule
\textbf{Total} & \textbf{130} \\
\bottomrule
\end{tabular}
\fonte{\texttt{docs/renders/mesh\_report.txt} (2026).}
\end{table}

\section{Objetos com arestas de borda}

Apenas tr\^es objetos apresentam arestas de borda, todos eles emblemas de texto:
\texttt{GEO\_Badge\_}\allowbreak\texttt{LOTUS} (1.240),
\texttt{GEO\_Badge\_}\allowbreak\texttt{111S} (536) e, adicionalmente, as
bordas de contorno dos demais glifos. Nenhum objeto de superf\'icie curva apresenta aresta de
borda, o que confirma que os corpos s\~ao fechados e as tampas operam conforme projetado.

\section{Objetos com tri\^angulos}

Os tri\^angulos concentram-se em tr\^es origens: os glifos dos emblemas
(\texttt{GEO\_Badge\_}\allowbreak\texttt{LOTUS} com 604 e
\texttt{GEO\_Badge\_}\allowbreak\texttt{111S} com 252), os cantos de chanfro
de primitivas de caixa e os elementos \'opticos de revolu{\c{c}}\~ao. Nenhum tri\^angulo
ocorre em superf\'icie curva principal do casco ou da cabine.
